\section{Verified quantized cached inference}
\label{sec:scope}

A Lean-checked theorem connects the exact 28,315-byte WebAssembly binary to the specified quantized cached-inference algorithm, including validation, termination, allocation, and buffer release for sessions of at most 128 tokens.  The model has GPT-2's 124,439,808 parameters, twelve transformer blocks, hidden width 768, twelve attention heads, feed-forward width 3,072, and vocabulary size 50,257.  The checkpoint stores 1,024 position embeddings.  The deployed recurrence accepts positions zero through 127.  We use the existing checkpoint, tokenizer, prompts, serial binary32 implementation, and reference harness from the earlier cached-inference work~\cite{gpt2,fp32report,comprehensivereport}.

The new computation quantizes learned projection weights and their activation inputs to integers in $[-127,127]$.  Each group of 64 activation coordinates has its own scale.  Every output channel has one weight scale across its input width.  Integer products accumulate in 32-bit words, and each group sum converts to binary32 before rescaling and addition.  Token embeddings share the quantized vocabulary matrix.  Position embeddings, normalization parameters, biases, attention scores and values, cached keys and values, GELU, softmax, residual additions, and output logits retain binary32 arithmetic.

The execution proof establishes exact agreement with the quantized algorithm for the frozen binary.  Experiments measure the algorithm's agreement with its independent reference and its differences from the binary32 model.  Numerical theorems bound those differences under explicit assumptions.  The evaluated forward bounds are too large to establish a token choice.  Certificates constructed from already observed logit pairs establish some individual choices.  These conclusions have different premises and evidence, which Sections~\ref{sec:proof}, \ref{sec:accuracy}, and~\ref{sec:numerical} specify.

The report fixes the implementation and evidence at LeanExe revision \longcode{c655d35b5011c1703dfd22bcceaec4e5bee17088}.  Quantized command-line inference is opt-in through \code{tools/gpt2 --quantized}.  The binary32 path remains the default.  The focused quantized source and exact-binary package checks have passed.  Repository-wide release work remains open at this revision: current limit/market matcher proof repairs, their preserved artifact packages, aggregate source and artifact checks, the complete execution and conformance tests, and a matching clean-checkout receipt.  The reported GPT results therefore identify a checked component and recorded experiments within an unfinished repository release.

\subsection{Relationship to quantization and verification work}

Integer quantization with explicit scale and rounding rules follows established inference practice.  Jacob et al. describe integer-only inference and training designed to preserve accuracy~\cite{jacob}.  The present implementation uses post-training quantization and binary32 rescaling and nonlinear operations.  Its deployed mathematical object is a mixed integer/binary32 recurrence with a fixed evaluation order.

Activation outliers have motivated specialized transformer quantizers.  LLM.int8() combines vector-wise quantization with a separate higher-precision computation for outlier dimensions~\cite{llmint8}.  SmoothQuant transfers activation scaling difficulty into weights through an equivalent transformation before quantization~\cite{smoothquant}.  Our experiment instead changes the activation scale granularity to 64-coordinate groups while preserving eight-bit products and the exported coefficient payload.  The comparison in Section~\ref{sec:accuracy} measures this particular change on GPT-2.  The benchmark compares two LeanExe WebAssembly implementations on one host, so it supplies no runtime comparison with those systems.

QEBVerif uses differential reachability analysis and mixed-integer linear programming to verify bounds between neural networks and their quantized counterparts~\cite{qebverif}.  Our numerical work proves compositional error formulas in Lean for the integer and binary32 operations of a cached transformer.  The binary proof additionally covers the deployed instructions and memory operations.  The present forward evaluator's failure to establish useful margins is a concrete limitation of its error formulas.  The earlier LeanExe reports supply the binary32 session, floating-point semantics, and proof organization reused here~\cite{fp32report,comprehensivereport}.
